\chapter{МЕТОДЫ И ПОСТАНОВКА ЗАДАЧИ}

\section{Формализация задачи}

В подразделе 1.5 идентифицирован пробел в существующих решениях, относящийся к задаче формирования согласованных рекомендательных выдач для множества пользователей одновременно при жёстких ограничениях вместимости. Настоящий подраздел представляет формальную постановку задачи, опирающуюся на синтез двух концептуальных рамок, рассмотренных в подразделах 1.3.2 и 1.3.4: теории игр с перегрузкой (congestion games) и марковского процесса принятия решений с ограничениями (Constrained MDP). Первая рамка обеспечивает структурное описание ситуации, в которой коллективное распределение пользователей формируется как результат децентрализованных индивидуальных выборов под влиянием информационного сигнала; вторая — операциональный аппарат для построения управляющей политики, обучаемой методами обучения с подкреплением.

\subsection{Объекты и переменные задачи}

Рассматривается дискретное множество ресурсов, обозначаемых индексом из множества $J = \{1, \ldots, m\}$, и дискретное множество пользователей с индексами из множества $I = \{1, \ldots, n\}$. Каждый ресурс $j$ характеризуется вместимостью $c_j$, представляющей предельное число пользователей, для которых обслуживание ресурсом возможно без деградации качества. Каждый пользователь $i$ характеризуется профилем $p_i$ — описанием его интересов в форме, поддерживающей численное вычисление меры близости с описаниями ресурсов; в работе профиль представлен векторным эмбеддингом текстового описания интересов, согласованным по размерности с эмбеддингами описаний ресурсов.

Время рассматривается дискретно: эпизод взаимодействия пользователей с системой разбивается на тайм-слоты $T = \{1, \ldots, \tau\}$, в каждом из которых пользователю предъявляется рекомендательная подсказка и наблюдается фактический выбор. В пределах одного тайм-слота пользователь может быть сопоставлен не более чем с одним ресурсом; полный план взаимодействия пользователя с системой за эпизод формируется как последовательность сопоставлений по тайм-слотам.

\subsection{Управляющая политика}

Управляющая сторона реализует политику рекомендаций $\pi$, отображающую состояние системы в распределение возможных рекомендательных выдач. На каждом тайм-слоте $t$ для каждого пользователя $i$ политика формирует ранжированный список $L_{i,t}$ из $K$ ресурсов, отсортированных в порядке убывания предсказанной полезности. Список $L_{i,t}$ является непосредственным управляющим воздействием на пользователя; пользователь самостоятельно принимает решение о выборе ресурса, опираясь на представленный список и собственные предпочтения.

Существенным элементом постановки является зависимость рекомендации одного пользователя от рекомендаций, выдаваемых остальным пользователям в тот же тайм-слот. Формирование списка $L_{i,t}$ изолированно от $\{L_{j,t}\}_{j \neq i}$ воспроизводит классическую постановку индивидуальных рекомендательных систем и приводит, как показано в подразделе 1.2, к коллективным эффектам перегрузки. В рамках предлагаемой формализации политика $\pi$ принимает на вход агрегированное состояние системы, включающее текущую загрузку ресурсов, что обеспечивает информационную связность рекомендаций между пользователями.

\subsection{Релевантность как функция близости}

Полезность сопоставления пользователя $i$ с ресурсом $j$ определяется через меру семантической близости между профилем пользователя и описанием ресурса. Используется косинусная близость векторных эмбеддингов:

\[
r_{ij} = \cos(p_i, e_j) = \frac{p_i \cdot e_j}{\|p_i\| \cdot \|e_j\|},
\]

где $p_i$ — эмбеддинг профиля пользователя, $e_j$ — эмбеддинг описания ресурса. Выбор косинусной близости в качестве функции релевантности обусловлен её интерпретируемостью, низкой вычислительной стоимостью и широким распространением в качестве базового сигнала ранжирования в литературе рекомендательных систем.

Чисто семантическая близость в роли сигнала релевантности порождает известную проблему смыслового сужения выдачи (filter bubble): пользователь систематически получает рекомендации только тематически близких ресурсов. В работе эта проблема решается не на уровне определения релевантности, а на уровне сравниваемых политик: часть из них (в частности, MMR и LLM-ranker) предусматривают механизмы преодоления смыслового сужения, что позволяет провести систематическое сравнение в этом отношении.

\subsection{Отклик пользователя}

Ответ пользователя $i$ на рекомендательный список $L_{i,t}$ моделируется как случайная величина $y_{i,t}$, принимающая значения из множества $J \cup \{\varnothing\}$, где значение $\varnothing$ соответствует отказу от выбора в данном тайм-слоте. Распределение $y_{i,t}$ задаётся условной вероятностной функцией выбора $P(y_{i,t} \mid L_{i,t}, p_i, s_t)$, зависящей от рекомендательного списка $L_{i,t}$, профиля пользователя $p_i$ и текущего состояния системы $s_t$ (в первую очередь — текущей заполненности ресурсов). Вероятностный характер отклика отражает фундаментальное ограничение управляющего рычага: рекомендация формирует пользовательскую перцепцию доступных вариантов, но не предопределяет фактического выбора. Базовая функция выбора — мультиномиальная логит-модель с температурой: вероятность выбора ресурса пропорциональна экспоненте релевантности с учётом штрафа за заполненность, отказ представлен отдельной альтернативой с базовой вероятностью; точные параметры приводятся в главе 3.

Все компоненты совокупной награды рассчитываются по реализациям случайных величин $\{y_{i,t}\}$, а не по содержанию рекомендательных списков $\{L_{i,t}\}$: компонент пользовательской полезности агрегирует релевантности $r_{i, y_{i,t}}$ для тех пользователей и тайм-слотов, в которых $y_{i,t} \neq \varnothing$; мера переполнения и дисперсия заполненности рассчитываются по фактическому распределению пользователей по ресурсам, индуцированному реализациями $\{y_{i,t}\}$. Состояние $s_t$, наблюдаемое политикой на каждом шаге, аналогично формируется по реализациям $y_{i,t}$ для пользователей, обработанных ранее в текущем тайм-слоте. Тем самым проигнорированная пользователем рекомендация не приносит положительного вклада в полезность и не уменьшает наблюдаемой загрузки альтернативных ресурсов, что естественным образом штрафует политику за выдачу рекомендаций, не согласованных с предпочтениями пользователя.

\subsection{Целевая функция управляющей стороны}

Целью управляющей стороны является максимизация совокупного результата эпизода, агрегирующего три компонента: суммарную пользовательскую полезность выбранных ресурсов, минимизацию переполнений вместимости, минимизацию неравномерности загрузки ресурсов. Совокупная награда эпизода формируется как взвешенная сумма указанных компонентов:

\[
R = \alpha \cdot R_{\text{rel}} - \beta \cdot R_{\text{over}} - \gamma \cdot R_{\text{var}},
\]

где $R_{\text{rel}}$ — суммарная релевантность фактически выбранных пользователями ресурсов; $R_{\text{over}}$ — мера переполнения, агрегированная по тайм-слотам как сумма превышений вместимостей в долях от соответствующих лимитов; $R_{\text{var}}$ — дисперсия заполненности ресурсов, агрегированная по тайм-слотам; $\alpha, \beta, \gamma$ — неотрицательные весовые коэффициенты, определяющие относительную приоритетность компонентов. Конкретные значения коэффициентов определяются на этапе настройки и приводятся в главе 4.

\subsection{Постановка как марковский процесс принятия решений с ограничениями}

Описанная задача формализуется как марковский процесс принятия решений с ограничениями (Constrained MDP) в смысле, рассмотренном в подразделе 1.3.4. На каждом шаге $t$ управляющий агент наблюдает состояние $s_t$, включающее текущего обслуживаемого пользователя, его профиль, агрегированные характеристики уже обработанных пользователей и текущую загрузку ресурсов; выбирает действие $a_t = L_{i,t}$ — ранжированный список из $K$ ресурсов; получает мгновенную награду $r_t$, определяемую вкладом тайм-слота в совокупную награду $R$ согласно приведённой выше формуле; и переходит в новое состояние $s_{t+1}$, отражающее последствия выбора пользователя. Целью является нахождение стационарной политики $\pi^*$, максимизирующей ожидание совокупной награды по эпизоду.

\subsection{Связь с играми с перегрузкой}

Структурная природа задачи соответствует концепции игр с перегрузкой (congestion games) в смысле подраздела 1.3.2. Пользователи представляют автономных агентов, выбирающих ресурсы из предложенных вариантов; полезность ресурса для конкретного пользователя содержит компонент, зависящий от числа других пользователей, выбравших тот же ресурс (через ограничение вместимости и связанные с ним эффекты). В отсутствие управляющего вмешательства система движется к равновесию эгоистичных агентов, отдалённому от системного оптимума на величину, измеряемую показателем «цена анархии». Управляющая политика $\pi$ выступает в роли информационного сигнала, формирующего издержки альтернативных выборов в восприятии каждого пользователя и тем самым смещающего равновесие в направлении системно-оптимального распределения.

Предлагаемая постановка, таким образом, объединяет аппарат двух рамок: постановка задачи как Constrained MDP обеспечивает операциональный механизм построения и обучения управляющей политики, а трактовка через игру с перегрузкой — концептуальное обоснование необходимости совместной оптимизации рекомендаций между пользователями и теоретическую интерпретацию роли управляющей политики как механизма коррекции децентрализованного равновесия.

\section{Требования к решению}

Постановка задачи, сформулированная в подразделе 2.1, и анализ ограничений существующих подходов, выполненный в подразделе 1.5, позволяют сформулировать требования, которым должно удовлетворять разрабатываемое решение. Требования разделяются на три группы: функциональные требования, описывающие необходимые возможности системы; нефункциональные требования, описывающие её качественные характеристики; ограничения предметной области, определяющие условия, в которых решение должно работать.

\subsection{Функциональные требования}

К числу функциональных требований к разрабатываемому решению относятся:

\begin{enumerate}
  \item \textbf{Приём входных данных.} Система должна принимать на вход множество ресурсов с описаниями и вместимостями, а также множество пользователей с профилями интересов. Представление профилей и описаний должно поддерживать численное вычисление меры близости.
\end{enumerate}

\begin{enumerate}
  \item \textbf{Формирование персональных рекомендаций.} Для каждого пользователя в каждом тайм-слоте система должна формировать ранжированный список из $K$ ресурсов, упорядоченных по предсказанной полезности с учётом текущего состояния системы.
\end{enumerate}

\begin{enumerate}
  \item \textbf{Учёт ограничений вместимости.} Формирование рекомендаций должно учитывать вместимости ресурсов и текущую их заполненность; политика рекомендаций должна избегать систематического направления пользовательского потока к ресурсам, для которых превышение вместимости становится вероятным.
\end{enumerate}

\begin{enumerate}
  \item \textbf{Согласованность рекомендаций между пользователями.} Рекомендация, выдаваемая одному пользователю, должна формироваться с учётом рекомендаций, выданных и реализованных другими пользователями в текущем тайм-слоте.
\end{enumerate}

\begin{enumerate}
  \item \textbf{Обработка вероятностного отклика.} Система должна корректно обрабатывать стохастический характер пользовательского отклика, включая возможность выбора ресурса вне рекомендательного списка и отказа от выбора в тайм-слоте.
\end{enumerate}

\begin{enumerate}
  \item \textbf{Обновление состояния по фактическим выборам.} Состояние системы, наблюдаемое политикой, должно обновляться по реализованным пользовательским выборам, а не по выданным рекомендациям.
\end{enumerate}

\begin{enumerate}
  \item \textbf{Поддержка сравнения нескольких политик.} Система должна поддерживать прозрачное сравнение различных рекомендательных политик на одной и той же среде с использованием единого набора метрик.
\end{enumerate}

\subsection{Нефункциональные требования}

К числу нефункциональных требований относятся:

\begin{enumerate}
  \item \textbf{Воспроизводимость.} Прогон системы при фиксированных входных данных и параметрах среды должен давать одинаковый результат, что обеспечивается фиксацией начальных значений генераторов случайных чисел и явной спецификацией всех параметров эксперимента.
\end{enumerate}

\begin{enumerate}
  \item \textbf{Эффективность обучения.} Обучение и сравнение политик должны быть осуществимы в разумное время на доступном исследовательском оборудовании уровня одной графической вычислительной карты потребительского класса.
\end{enumerate}

\begin{enumerate}
  \item \textbf{Расширяемость.} Архитектура решения должна поддерживать добавление новых рекомендательных политик и новых функций отклика пользователей без переработки симуляционной среды и компонентов оценки качества.
\end{enumerate}

\begin{enumerate}
  \item \textbf{Прозрачность.} Все компоненты совокупной награды и итоговые метрики должны допускать содержательную интерпретацию в терминах задачи, что необходимо как для отладки, так и для последующего обоснования результатов работы.
\end{enumerate}

\begin{enumerate}
  \item \textbf{Низкая стоимость обучения.} Затраты на обучение и сравнение политик не должны существенно зависеть от стоимости обращения к большим языковым моделям. Использование больших языковых моделей допускается на этапах, требующих однократных или редких обращений, но не на каждом шаге обучения.
\end{enumerate}

\begin{enumerate}
  \item \textbf{Внешняя валидность результатов.} Качество результирующих политик должно подтверждаться не только в среде, в которой проходило обучение, но и в независимой среде с альтернативной моделью пользовательского поведения, что обеспечивает контроль разрыва между симуляцией и реальным поведением пользователей.
\end{enumerate}

\subsection{Ограничения предметной области}

Дополнительно фиксируются ограничения, определяющие условия, в которых решение должно сохранять работоспособность:

\begin{enumerate}
  \item \textbf{Малый каталог ресурсов.} Число рассматриваемых ресурсов исчисляется десятками, что исключает применение методов, требующих больших каталогов для устойчивого обучения базовых моделей релевантности.
\end{enumerate}

\begin{enumerate}
  \item \textbf{Отсутствие журналов реальных пользовательских взаимодействий.} Решение должно обеспечивать обучение и сравнение политик в условиях отсутствия исторических данных о фактических пользовательских выборах.
\end{enumerate}

\begin{enumerate}
  \item \textbf{Ограниченный бюджет вычислений и обращений к большим языковым моделям.} Объём обучающих и оценочных вычислений ограничен ресурсами одного исследователя; объём обращений к большим языковым моделям — соответствующим бюджетом.
\end{enumerate}

Совокупность сформулированных требований формирует основу для описания предлагаемого метода в подразделе 2.3.

\section{Описание предлагаемого метода}

Предлагаемый метод объединяет три компонента: гибридную симуляционную среду для обучения и оценки политик при отсутствии исторических данных; набор сравниваемых рекомендательных политик различной природы — от случайной до обучаемой методами обучения с подкреплением и основанной на большой языковой модели; двухслойную схему валидации с операциональной проверкой в обучающей среде и независимой проверкой в среде с альтернативной моделью пользовательского поведения. Каждый компонент адресует подмножество требований из подраздела 2.2.

\subsection{Гибридная симуляционная среда}

Симуляционная среда, в которой осуществляется обучение и сравнение рекомендательных политик, имеет двухкомпонентную структуру.

\textbf{Компонент однократной генерации пользовательской популяции.} Перед началом обучения большая языковая модель формирует синтетическую популяцию пользователей в количестве, достаточном для статистически значимого сравнения политик. Каждый пользователь характеризуется текстовым профилем интересов, отражающим его предметные предпочтения, уровень квалификации, тематическую специфику и иные характеристики, влияющие на пользовательский выбор. Сгенерированные профили преобразуются в векторные эмбеддинги, согласованные по размерности с эмбеддингами описаний ресурсов. Однократный характер обращения к большой языковой модели существенно снижает совокупную стоимость экспериментов по сравнению со схемами, в которых обращение к большой языковой модели выполняется на каждом шаге обучения, при сохранении содержательного разнообразия пользовательской популяции. Использование большой языковой модели в роли генератора синтетических данных описано в подразделе 1.4.4 как одно из основных направлений применения языковых моделей в задачах рассматриваемого класса.

\textbf{Компонент параметрического пошагового симулятора.} Основной цикл обучения и оценки политик реализуется параметрическим симулятором, разыгрывающим на каждом шаге функцию выбора пользователей. Функция выбора реализуется мультиномиальной логит-моделью с температурой, описанной в подразделе 2.1, и зависит от релевантности предложенных ресурсов профилю пользователя, текущей заполненности ресурсов и базовой склонности пользователя к отказу от участия в тайм-слоте. Параметры функции выбора (температура, штраф за заполненность, базовая вероятность отказа) фиксируются на этапе настройки симулятора. Параметрический характер компонента обеспечивает высокую скорость симуляции, воспроизводимость результатов при фиксации начальных значений генератора случайных чисел и независимость стоимости обучения от стоимости обращений к большой языковой модели.

Гибридная структура среды отвечает требованию низкой стоимости обучения (нефункциональное требование 5) при сохранении содержательного разнообразия популяции, недостижимого в чисто параметрических симуляторах. Кроме того, она удовлетворяет требованиям воспроизводимости и эффективности (нефункциональные требования 1 и 2).

\subsection{Множество сравниваемых рекомендательных политик}

В рамках предлагаемого метода рассматривается шесть рекомендательных политик различной природы, обеспечивающих систематическое сравнение различных подходов к решению задачи. Каждая политика реализует отображение текущего состояния системы и профиля обслуживаемого пользователя в ранжированный список из $K$ ресурсов.

\textbf{Политика 1: Случайный выбор.} Формирует ранжированный список путём равномерно случайного выбора $K$ ресурсов из каталога без учёта профиля пользователя и состояния системы. Не предполагает обучения. Назначение в составе сравнения — установление нижней планки качества, относительно которой содержательность остальных политик должна быть очевидна; при отсутствии существенного превышения над случайной политикой возникают основания усомниться в адекватности постановки задачи или метрик.

\textbf{Политика 2: Косинусная близость без учёта ограничений.} Формирует список как первые $K$ ресурсов, упорядоченных по убыванию косинусной близости между эмбеддингом профиля пользователя и эмбеддингами описаний ресурсов. Соответствует классическому подходу контентной фильтрации в рекомендательных системах без учёта системных ограничений. Назначение — представление точки отсчёта для подходов, не учитывающих ограничения вместимости и взаимосвязанность рекомендаций между пользователями.

\textbf{Политика 3: Косинусная близость с учётом вместимости (rule-based).} Расширение второй политики: к косинусной релевантности применяется мультипликативный штраф, монотонно убывающий с ростом текущей заполненности соответствующего ресурса; ресурсы, заполненные выше заданного порога, исключаются из ранжированного списка. Не предполагает обучения и реализует управление через детерминированное правило. Назначение — представление эвристического подхода, учитывающего ограничения вместимости без обращения к обучаемым методам.

\textbf{Политика 4: Maximal Marginal Relevance.} Реализует алгоритм MMR, описанный в подразделе 1.4.1, с балансом релевантности и разнообразия выдачи через параметр $\lambda$. Не предполагает обучения. Назначение — представление классического академического подхода к преодолению проблемы смыслового сужения выдачи; обеспечивает контроль того, насколько добавление разнообразия влияет на загрузку различных ресурсов в условиях рассматриваемой задачи.

\textbf{Политика 5: Обучаемая политика на основе обучения с подкреплением.} Реализуется методом проксимальной оптимизации политики (PPO), описанным в подразделе 1.4.2. Обучается в параметрическом симуляторе на совокупной награде, определённой в подразделе 2.1. Состояние, наблюдаемое политикой, включает профиль текущего обслуживаемого пользователя, текущую заполненность ресурсов и агрегированные характеристики уже обработанных пользователей. Действие политики — ранжированный список из $K$ ресурсов. Назначение — представление обучаемого подхода, способного учитывать структуру ограничений и согласованность рекомендаций между пользователями за счёт обучения непосредственно в симуляторе. Эта политика — основной технический результат работы.

\textbf{Политика 6: LLM-ranker.} Реализует подход, в котором большая языковая модель выступает в роли ранжирующего компонента. Модель получает на вход профиль пользователя, описания ресурсов из числа $N$ ресурсов с наибольшей косинусной близостью к профилю и сводку текущей заполненности залов; на выходе формирует ранжированный список из $K$ ресурсов с учётом всех представленных факторов. Соответствует подходу, описанному в подразделе 1.4.4 как одно из основных направлений применения больших языковых моделей в задачах рассматриваемого класса. Назначение — сравнение с современным альтернативным методом, представляющим парадигму непосредственного использования больших языковых моделей в качестве рекомендательного компонента.

Совокупность шести политик покрывает основные классы подходов, рассмотренные в подразделах 1.3 и 1.4: случайный выбор как формальная нижняя граница; чистое контентное ранжирование как классический recsys-подход без ограничений; правиловое управление с учётом ограничений; классический алгоритм диверсификации; обучаемая политика в постановке Constrained MDP; современный подход на основе большой языковой модели. Подобный охват обеспечивает систематическое сравнение и позволяет позиционировать обучаемую политику относительно как простых эвристик, так и современных альтернатив.

\subsection{Двухслойная схема валидации}

Внешняя валидность результатов обеспечивается двухслойной схемой проверки, разделяющей этапы обучения и сравнения политик и этап независимой проверки результата.

\textbf{Слой 1: основное сравнение в параметрическом симуляторе.} Все шесть политик обучаются (для политик 5 и 6 — настраиваются) и сравниваются в гибридной симуляционной среде, описанной выше. На этом слое реализуется протокол сравнения по совокупности метрик, определяемых в подразделе 2.4, и выполняется выбор лучшей политики по совокупному критерию. Среда первого слоя характеризуется высокой скоростью, воспроизводимостью и низкой стоимостью.

\textbf{Слой 2: независимая валидация лучшей политики в LLM-агентной среде.} Политика, признанная лучшей по результатам первого слоя, проходит дополнительную проверку в среде, в которой пользователи моделируются не параметрической мультиномиальной логит-моделью, а автономными агентами на основе больших языковых моделей с собственной памятью и функцией выбора, как описано в подразделе 1.4.4 со ссылкой на работу Agent4Rec. Среда второго слоя характеризуется существенно большей вычислительной стоимостью, что ограничивает её применение этапом итоговой проверки. Сопоставление результатов первого и второго слоёв обеспечивает оценку устойчивости выбранной политики к изменению модели пользовательского поведения, тем самым контролируя разрыв между симуляцией и реальным поведением (sim-to-real gap).

Двухслойная схема непосредственно адресует требование внешней валидности результатов: разделение слоёв по стоимости обеспечивает совместимость требований внешней валидности и низкой стоимости обучения, которые в одноуровневой схеме противоречат друг другу.

\subsection{Соответствие требованиям}

Предлагаемый метод покрывает все требования, сформулированные в подразделе 2.2. Функциональные требования 1–7 обеспечиваются совокупностью симуляционной среды (требования 1, 5, 6) и сравниваемых политик (требования 2, 3, 4, 7), причём учёт ограничений вместимости и согласованность рекомендаций между пользователями реализуются на уровне состояния, наблюдаемого политикой. Нефункциональные требования удовлетворяются гибридной структурой среды (требования 1, 2, 5), архитектурой раздельно обучаемых политик (требование 3), интерпретируемой структурой совокупной награды (требование 4) и двухслойной схемой валидации (требование 6). Ограничения предметной области учитываются за счёт переноса основной нагрузки обучения в параметрический симулятор (ограничение 3), синтетической генерации пользовательской популяции (ограничение 2) и поддержки малого каталога ресурсов всеми компонентами метода (ограничение 1).

\section{Сравнительный анализ альтернатив}

Помимо предлагаемого метода, описанного в подразделе 2.3, существуют альтернативные методологические подходы, опирающиеся на другие концептуальные рамки из числа рассмотренных в подразделе 1.3 либо на иные конфигурации симуляционных сред. Настоящий подраздел проводит сравнительный анализ четырёх альтернатив с предлагаемым методом и обосновывает выбор последнего.

\subsection{Альтернатива 1. Чисто оптимизационный подход в рамках исследования операций}

Альтернативой предлагаемому методу является формулировка задачи как обобщённой задачи о назначениях (GAP) или многопериодной задачи о назначениях, рассмотренной в подразделе 1.3.1, с последующим её решением точными методами целочисленного программирования или приближёнными эвристиками. Подобный подход обеспечивает строгую математическую постановку и развитый алгоритмический арсенал.

Существенным препятствием для применения данной альтернативы в качестве основного метода является фундаментальное допущение, лежащее в основе классической постановки задач о назначениях: решение управляющей стороны полагается исполнимым, то есть приравнивается к фактическому распределению пользователей по ресурсам. В рассматриваемой задаче решение управляющей стороны имеет статус рекомендации, которой пользователь может следовать или не следовать. Прямое применение подхода исследования операций к задаче рекомендательного управления приводит к решениям, оптимальным в детерминированной модели, но потенциально неоптимальным с учётом вероятностного отклика пользователей.

Вместе с тем, оптимизационный подход сохраняет ценность в составе предлагаемого метода: решение детерминированной задачи о назначениях используется как теоретическая верхняя граница достижимой суммарной полезности при идеальном следовании пользователей рекомендациям, что обеспечивает контрольную точку для интерпретации результатов сравниваемых политик.

\subsection{Альтернатива 2. Чисто теоретико-игровой подход}

Альтернативой является описание задачи как игры с перегрузкой в рамках теоретико-игровой рамки подраздела 1.3.2 с последующим аналитическим вычислением равновесия Нэша или равновесия Уордропа, а также показателя «цена анархии» для системы. Подобный подход опирается на развитую теорию и обеспечивает концептуально чистое описание ситуации децентрализованного выбора автономных агентов.

Препятствием для применения данной альтернативы в качестве основного метода является описательный, а не предписывающий характер классической теоретико-игровой рамки. Аппарат равновесий отвечает на вопрос, какое распределение установится при заданных правилах взаимодействия и функциях издержек, однако не предоставляет алгоритма построения конкретного управляющего воздействия — рекомендательной выдачи, — способного сместить поведение пользователей в направлении системного оптимума. Переход от описания равновесий к проектированию управляющих воздействий требует обращения к смежным теоретическим направлениям, развитым преимущественно для целей теоретического анализа и не предоставляющим готовых вычислительных алгоритмов.

В предлагаемом методе теоретико-игровая рамка используется как структурное обоснование необходимости совместной оптимизации рекомендаций между пользователями (подраздел 2.1) и теоретическая интерпретация роли управляющей политики, но операциональным методом построения политики выступает аппарат марковского процесса принятия решений с ограничениями.

\subsection{Альтернатива 3. Прямое обучение в LLM-агентной среде}

Альтернативой является использование LLM-агентной симуляционной среды, описанной в подразделе 1.4.4, в качестве основной среды для обучения рекомендательной политики, без параметрического симуляционного слоя. Подобная среда обеспечивает максимальную содержательную глубину пользовательского поведения, моделируя каждого пользователя как автономного агента на основе большой языковой модели с собственной памятью и функцией выбора.

Препятствием для применения данной альтернативы является вычислительная стоимость. Обучение рекомендательной политики методами обучения с подкреплением требует значительного числа взаимодействий со средой — порядка нескольких миллионов шагов в типовых конфигурациях. В LLM-агентной среде каждый шаг включает множественные обращения к большой языковой модели, что превращает совокупный бюджет обучающего эксперимента в непрактично высокий при ограниченных вычислительных ресурсах.

В предлагаемом методе LLM-агентная среда применяется как слой 2 валидации, охватывающий только итоговую проверку лучшей политики на ограниченном числе сценариев, что укладывается в реалистичный экспериментальный бюджет.

\subsection{Альтернатива 4. Прямое применение классических методов рекомендательных систем с ограничениями}

Альтернативой является прямое применение методов рекомендательных систем с ограничениями, описанных в подразделе 1.3.3, — в первую очередь, постановки ранжирования с ограничениями Сингха и Йоахимса. Указанные методы обеспечивают формальный аппарат и алгоритмы оптимизации политики ранжирования при дополнительных ограничениях на агрегированную экспозицию ресурсов.

Препятствием для применения данной альтернативы является принципиальное отличие в типе формулируемого ограничения. Классические методы рекомендательных систем с ограничениями формулируют ограничения как математическое ожидание по распределению пользователей, что подразумевает выравнивание совокупной экспозиции ресурсов «в среднем» по большому потоку независимых выдач. В рассматриваемой постановке ограничения вместимости должны соблюдаться в каждой конкретной коллективной выдаче, формируемой одновременно для фиксированного множества пользователей. Адаптация классических методов под подобный режим работы требует существенной методологической переработки, по объёму сопоставимой с разработкой нового метода.

Вместе с тем, методы данного направления находят применение в составе сравниваемых политик предлагаемого метода — в частности, политика 4 (Maximal Marginal Relevance) реализует один из канонических алгоритмов переранжирования, описанных в подразделе 1.4.1.

\subsection{Сводное обоснование выбора}

Сравнительный анализ показывает, что ни одна из рассмотренных альтернатив, взятая отдельно, не закрывает совокупности требований, сформулированных в подразделе 2.2. Подход исследования операций (альтернатива 1) не учитывает вероятностного отклика; теоретико-игровой подход (альтернатива 2) описателен, а не предписывающ; прямое обучение в LLM-агентной среде (альтернатива 3) непрактично по стоимости; прямое применение методов рекомендательных систем с ограничениями (альтернатива 4) опирается на формализм амортизированных, а не одноразовых ограничений.

Предлагаемый метод — синтез элементов нескольких рамок: операциональная постановка взята из аппарата марковского процесса принятия решений с ограничениями (подраздел 1.3.4), концептуальное обоснование совместной оптимизации опирается на теорию игр с перегрузкой (подраздел 1.3.2), сравниваемые политики покрывают разные классы подходов из подразделов 1.3.3, 1.4.1 и 1.4.4, двухслойная схема валидации согласует требования низкой стоимости обучения и внешней валидности результатов. Такая конфигурация удовлетворяет всем требованиям одновременно, что не достижимо ни одной из альтернатив в её прямой форме.

\section{Дизайн экспериментов и план валидации}

Настоящий подраздел определяет дизайн экспериментальной части работы, направленной на сравнительную оценку шести политик, описанных в подразделе 2.3, и независимую валидацию лучшей политики во втором слое симуляционной среды. Экспериментальная часть включает четыре составляющие: сравнение политик по основным метрикам качества, исследования влияния отдельных компонентов метода (ablation studies), проверку устойчивости результатов к вариации параметров среды и итоговую валидацию в LLM-агентной среде. Все экспериментальные процедуры опираются на единый набор метрик, описанный в первой части подраздела.

\subsection{Основные метрики оценки}

Сравнение рекомендательных политик опирается на четыре метрики, каждая из которых отражает отдельный аспект качества решения.

\textbf{Overflow rate.} Относительная мера переполнений ресурсов. Для каждого ресурса в каждом тайм-слоте вычисляется превышение фактического числа пользователей над вместимостью, отнесённое к вместимости; итоговое значение — среднее этих превышений по всем ресурсам, тайм-слотам и эпизодам. Метрика принимает значения от нуля (вместимости соблюдаются строго) до положительных значений, растущих с увеличением степени переполнения. Отражает соблюдение жёстких ограничений вместимости в каждой одноразовой выдаче.

\textbf{Hall utilization variance.} Дисперсия относительной заполненности ресурсов в пределах тайм-слота, усреднённая по тайм-слотам и эпизодам. Низкие значения метрики соответствуют равномерному распределению аудитории между ресурсами; высокие значения — ситуации одновременной перегрузки части ресурсов и недозагрузки других. Отражает согласованность рекомендаций, выдаваемых разным пользователям одновременно.

\textbf{Mean user utility.} Средняя релевантность ресурсов, фактически выбранных пользователями в ходе эпизода. Релевантность измеряется как косинусная близость профиля пользователя и описания выбранного ресурса (подраздел 2.1); пользователи, отказавшиеся от выбора в тайм-слоте, вносят нулевой вклад в среднее. Метрика отражает суммарную пользовательскую полезность, обеспечиваемую политикой.

\textbf{Gini coefficient по ресурсам.} Коэффициент Джини, рассчитанный по совокупному распределению аудитории между ресурсами за эпизод, отражающий неравенство получаемого ресурсами внимания. Близкие к нулю значения соответствуют равномерному распределению аудитории; близкие к единице — ситуации, в которой большая часть пользователей сосредоточена на малом числе ресурсов. Метрика обеспечивает оценку справедливости со стороны поставщиков ресурсов в смысле, рассмотренном в подразделе 1.3.3.

Совокупность четырёх метрик покрывает три аспекта качества: соблюдение системных ограничений (Overflow rate, Hall utilization variance), пользовательская полезность (Mean user utility), справедливость распределения внимания между ресурсами (Gini coefficient). Раздельное представление метрик обеспечивает возможность анализа компромиссов, возникающих между конкурирующими аспектами качества.

\subsection{Протокол сравнения политик}

Сравнение шести политик из подраздела 2.3 проводится в первом слое симуляционной среды по следующему протоколу.

В качестве экспериментальной обстановки фиксируется конкретная конфигурация среды: множество ресурсов с заданными описаниями и вместимостями, синтетическая популяция пользователей, сформированная однократной LLM-генерацией, расписание тайм-слотов. Все шесть политик прогоняются в идентичной среде; различие между прогонами определяется только начальным значением генератора случайных чисел, управляющего стохастикой функции выбора.

Каждая политика оценивается на множестве независимых эпизодов количества $n_{\text{ep}}$, отличающихся начальным значением случайных чисел. Для обучаемых политик (политика 5) обучение выполняется отдельно с использованием фиксированной части эпизодов, оценка — на отдельной отложенной выборке, что обеспечивает корректное разделение обучающих и оценочных данных. Для каждой метрики вычисляются среднее и доверительный интервал по эпизодам.

Сравнение результатов проводится по двум критериям. Первый критерий — раздельное сравнение по каждой из четырёх метрик; политики ранжируются по каждой метрике с указанием статистической значимости различий. Второй критерий — совокупное сравнение по взвешенной сумме компонентов совокупной награды, определённой в подразделе 2.1; этот критерий используется для выбора лучшей политики при наличии компромиссов между метриками.

\subsection{Исследование влияния отдельных компонентов метода (ablation studies)}

С целью оценки относительного вклада отдельных элементов предлагаемого метода в итоговое качество проводится серия исследований, в каждом из которых отключается или заменяется один из компонентов, после чего повторно проводится сравнение политик.

\textbf{Влияние LLM-генерации синтетической популяции.} Проводится сравнение результатов основного эксперимента с результатами эксперимента, в котором профили пользователей формируются равномерно случайно из пространства эмбеддингов либо отбираются случайно из заданного фиксированного набора. Цель — оценить, в какой мере содержательное разнообразие популяции, обеспечиваемое большой языковой моделью, сказывается на итоговых результатах сравнения политик.

\textbf{Влияние компонентов совокупной награды.} Проводится повторное обучение политики 5 (обучаемой методом проксимальной оптимизации политики) в трёх режимах: с полной совокупной наградой; с отключённым штрафом за переполнение ($\beta = 0$); с отключённым штрафом за дисперсию заполненности ($\gamma = 0$). Цель — определить относительный вклад каждой компоненты награды в результирующее поведение политики.

\textbf{Влияние информационной связности состояния.} Проводится повторное обучение политики 5 в режиме, в котором состояние, наблюдаемое политикой, не включает информацию о текущей заполненности ресурсов. Цель — эмпирически подтвердить значимость информационной связности рекомендаций между пользователями для решения задачи; ожидается, что отключение данной составляющей состояния приводит к существенной деградации метрик соблюдения вместимости.

\subsection{Проверка устойчивости результатов}

Устойчивость результатов сравнения политик к вариации параметров среды проверяется серией экспериментов, в которых варьируется один из параметров среды при фиксации остальных:

\begin{itemize}
  \item Объём пользовательской популяции — варьируется в диапазоне нескольких порядков для оценки масштабируемости методов;
  \item Распределение вместимостей ресурсов — рассматриваются равномерное и существенно неравномерное распределения;
  \item Параметр температуры функции выбора в симуляторе — варьируется в диапазоне, соответствующем переходу от детерминированного следования рекомендациям до близкого к равномерному выбору;
  \item Размер каталога ресурсов — варьируется в пределах, оставляющих каталог в режиме малого размера в смысле подраздела 2.2.
\end{itemize}

Для каждой комбинации параметров повторяется протокол сравнения политик. Согласованность ранжирования политик при различных значениях параметров рассматривается как индикатор устойчивости предлагаемого метода.

\subsection{Итоговая валидация во втором слое}

Лучшая политика по результатам первого слоя проходит независимую валидацию во втором слое симуляционной среды на основе LLM-агентов, описанной в подразделе 2.3. Протокол валидации включает прогон политики на ограниченном числе эпизодов в среде, в которой пользователи моделируются автономными агентами на основе больших языковых моделей с собственной памятью и функцией выбора.

В рамках второго слоя рассчитываются те же четыре метрики, что и в первом, и проводится их сопоставление с результатами первого слоя. Отдельное внимание уделяется качественному анализу расхождений: сценариям, в которых результаты двух слоёв существенно различаются. Подобный анализ обеспечивает оценку разрыва между симуляцией и реальным поведением (sim-to-real gap) и формирует материал для обсуждения внешней валидности результатов работы.

\subsection{Критерии успешности работы}

Совокупность экспериментов считается подтверждающей предлагаемый метод при выполнении трёх условий:

\begin{enumerate}
  \item Обучаемая политика (политика 5) демонстрирует статистически значимое превосходство над политиками 1–4 по совокупному критерию и по большинству отдельных метрик в первом слое симуляционной среды.
  \item Обучаемая политика 5 показывает результаты, сопоставимые с результатами политики 6 (LLM-ranker), либо превосходит её по совокупному критерию при существенно меньшей вычислительной стоимости отдельной рекомендации.
  \item Итоговая валидация во втором слое качественно подтверждает результаты первого слоя; наблюдаемые расхождения метрик находятся в пределах, не отменяющих основных выводов о ранжировании политик.
\end{enumerate}

Невыполнение одного или нескольких условий рассматривается как самостоятельный результат работы, подлежащий обсуждению в главе 4 в части анализа полученных результатов и формулировки направлений дальнейших исследований.

\section{Выводы по главе}

В главе сформулирована постановка задачи и описан предлагаемый методологический подход к её решению.

В подразделе 2.1 проведена формализация задачи на основе синтеза двух концептуальных рамок, рассмотренных в главе 1: марковского процесса принятия решений с ограничениями и игр с перегрузкой. Введены формальные обозначения объектов задачи, описана структура управляющей политики, определена функция релевантности на основе косинусной близости эмбеддингов профилей и описаний ресурсов, формализован вероятностный отклик пользователей через случайные величины фактических выборов. Целевая функция управляющей стороны представлена как взвешенная сумма пользовательской полезности, штрафа за переполнение вместимостей и штрафа за неравномерность загрузки ресурсов.

В подразделе 2.2 сформулированы требования, которым должно удовлетворять разрабатываемое решение. Требования разделены на три группы: семь функциональных требований к возможностям системы, шесть нефункциональных требований к её качественным характеристикам и три ограничения предметной области.

В подразделе 2.3 описан предлагаемый метод, объединяющий три компонента: гибридную симуляционную среду с однократной LLM-генерацией пользовательской популяции и параметрическим пошаговым симулятором; множество из шести сравниваемых рекомендательных политик различной природы; двухслойную схему валидации с основным сравнением в параметрическом симуляторе и независимой проверкой лучшей политики в среде на основе LLM-агентов. Показано, что предложенная конфигурация метода удовлетворяет всем требованиям, сформулированным в подразделе 2.2.

В подразделе 2.4 проведён сравнительный анализ четырёх альтернативных методологических подходов: чисто оптимизационного подхода в рамках исследования операций; чисто теоретико-игрового подхода; прямого обучения политики в LLM-агентной среде; прямого применения классических методов рекомендательных систем с ограничениями. Установлено, что ни одна из альтернатив, взятая отдельно, не закрывает совокупности требований, что обосновывает выбор предлагаемого метода как систематического синтеза элементов нескольких рамок.

В подразделе 2.5 определён дизайн экспериментальной части работы, включающий четыре составляющие: сравнение политик по основным метрикам качества (Overflow rate, Hall utilization variance, Mean user utility, Gini coefficient); исследования влияния отдельных компонентов метода (LLM-генерация персон, компоненты совокупной награды, информационная связность состояния); проверку устойчивости результатов к вариации параметров среды; итоговую валидацию лучшей политики во втором слое симуляционной среды. Сформулированы три критерия успешности работы.

На основании сформулированной постановки задачи, требований и описания метода в главе 3 представляются конкретные технические решения по реализации симуляционной среды и сравниваемых политик, а также описывается ход проведённых экспериментов и полученные результаты.
